% The appendix, in 5 groups. Each group is a section and each item a subsection,
% so the numbering stays within A to E whatever the appendix grows to.

\section{Setup, terms and attacks}
\label{app:setup}

\subsection{Terms and metrics}
\label{app:glossary}

\paragraph{Terms of the method.}
A sweep runs one placement on one backdoored model over its rate ladder and stores the per-pass predictions, so that every rule and metric below is read from the same stored passes. A probe is a placement together with the rate the rule assigns it on a given model, the unit a defender deploys. The basis is the list of placements we swept on every model, listed in \cref{app:basis}. A band restricts a placement to a range of blocks, and a placement without a band acts in every block. The benign reference of a dataset is a model trained with the same recipe and no poisoning, probed with the same triggers, which every mechanism experiment uses as its control.

\paragraph{Terms of the backdoor literature.}
We use the vocabulary of the backdoor literature. The poisoning rate is the share of the training set the attacker modifies. An attack is dirty-label when it changes the labels of the poisoned images and clean-label when it does not. It is all-to-one when every trigger maps to one target class and all-to-all when each source class maps to its own target. It is source-specific when only chosen source classes are poisoned. The attack success rate (ASR) is the share of triggered non-target test images classified as the target, and clean accuracy (CA) is the accuracy of the backdoored model on clean test images, whose drop against the benign model of the same dataset measures the attack's cost. For the detector, a positive is a triggered input. The true-positive rate (TPR, the detector's sensitivity) is the share of triggered inputs flagged, the false-positive rate (FPR, one minus its specificity) is the share of clean inputs flagged and AUROC is the area under the curve of TPR against FPR over every threshold. We write TPR@FPR for the TPR read at a fixed FPR budget, and the budget is set on the clean validation set.



\subsection{The attacks}
\label{app:attacks}

BadNets stamps a small fixed patch onto an image and swaps its label to a chosen target, one of the earliest data-poisoning backdoors for deep networks \cite{gu2017badnets}. Our patch is a 3 by 3 checkerboard, alternating black and white squares starting from white, stamped into the bottom right corner, the convention BackdoorBench uses on 32 by 32 inputs. The default run is dirty-label all-to-one at target class 0, and we separately train an all-to-all variant that maps each source class to the next class, wrapping the last class back to class 0, for the multiclass comparisons elsewhere in the paper. BadNets implants on every dataset and rate we test, with attack success no lower than \AttackBadnetATwooAsrMin{} even at \PanelRateLow{} poisoning. A control run trained at target class 1 on GTSRB, meant to check whether shifting the target confounds the clean-label results, collapsed instead to a single-class predictor mid-training, a training divergence covered in the protocol paragraph below rather than a property of this trigger.

Blend alpha-blends a fixed pattern over the whole image rather than localizing to a patch \cite{chen2017targeted}. We replace the paper's Hello Kitty image with a pattern of independent uniform pixels drawn once from a seeded generator, so the package ships no external image. We blend it at a ratio of \AttackBlendAlpha{} against the clean image. Labeling is dirty-label all-to-one at target class 0. Blend implants at every rate on every dataset, with attack success between \AttackBlendAsrMin{} and \AttackBlendAsrMax{}. Its own GTSRB target-1 control run collapsed in the same way as the BadNets one above.

SIG adds a horizontal sinusoidal signal across the image rather than a spatial patch and keeps every label unchanged, so only images that already carry the target label are perturbed \cite{barni2019sig}. We use Barni et al.'s own amplitude of \AttackSigAmplitude{} in 0-to-1 pixel units, at a frequency of \AttackSigFrequency{} cycles across the image width, clamped back into the valid pixel range after adding the signal. The label mode is clean-label with a single target class, class 0 on CIFAR-10, CIFAR-100 and Tiny ImageNet and class 1 on GTSRB, where class 0 is too small to reach a usable rate. SIG only clears the attack success bar on CIFAR-10 at 10\% poisoning. It never clears on CIFAR-100, GTSRB or Tiny ImageNet at any rate we could reach given the clean-label ceiling, staying between \AttackSigAsrNeverClearingMin{} and \AttackSigAsrNeverClearingMax{} there, which is why the results in the main text treat it as one of the hard attacks.

WaNet moves where pixels are sampled from rather than changing their values, warping the image through a smooth field generated from a small control grid \cite{nguyen2021wanet}. Our control grid is \AttackWanetControlGridSize{} by \AttackWanetControlGridSize{} points and the warp strength is \AttackWanetStrength{}. The field is drawn from seed 0 and shared by every poisoned image. Training also uses WaNet's noise mode, cover images warped by an independent random field per sample with their label kept, at \AttackWanetCoverRatio{} times the poisoning rate, following BackdoorBench's cross-ratio of 2 and the original paper. Labeling is dirty-label all-to-one at target class 0. WaNet never clears the attack success bar at 1\% poisoning on any dataset, reaching only \AttackWanetAsrLowRateMin{} to \AttackWanetAsrLowRateMax{} there. It clears at 5\% and 10\% on CIFAR-10 and Tiny ImageNet and only at 10\% on GTSRB. It never clears on CIFAR-100, where the best reading is \AttackWanetAsrNeverClearingMax{}.

LF adds a trigger whose energy sits in low spatial frequencies \cite{zeng2021lf}. Zeng et al. optimize it against a pretrained classifier and low-pass filter every step. We use a simpler variant that low-pass filters a fixed noise pattern in the Fourier domain and rescales it to span the full pixel range. Our filter keeps frequencies within a radius of \AttackLfCutoff{} around the spectrum center and adds the resulting pattern at a strength of \AttackLfStrength{}, clamped back into the valid pixel range. Labeling is dirty-label all-to-one at target class 0. LF implants at every rate on every dataset, with attack success between \AttackLfAsrMin{} and \AttackLfAsrMax{}, which puts it alongside BadNets, Blend and BPP among the attacks that implant most reliably.

Label-Consistent stamps a 3 by 3 black-and-white checkerboard, the pattern our BadNets also uses, into all four corners of a target-class image rather than one, and keeps that image's own label, so the trigger alone has to carry the backdoor \cite{turner2019label}. The patch-only version does not implant, because the network can still read the class off the untouched picture, so we follow Turner et al. and first perturb each base image with an untargeted L-infinity PGD attack \cite{madry2018pgd}, against a clean surrogate where they use an adversarially trained one, \AttackLcPgdSteps{} steps at a step size of \AttackLcPgdStepFactor{} times epsilon over the step count, so that the image stops supporting its own label. We swept \LcEpsilonBudgets{} epsilon budgets at each dataset's highest reachable clean-label rate and kept the smallest value that clears both the attack success bar and the clean-accuracy tolerance. The rule picks \LcEpsilonChosen{}, so \LcDatasetsWithEpsilon{} dataset passes it (\cref{tab:clean-label-lc-epsilon}). CIFAR-10 implants but loses more clean accuracy than the tolerance allows, and CIFAR-100 and Tiny ImageNet never reach the attack success bar at any epsilon we tried. Of the \LcGtsrbPilotRuns{} GTSRB pilot runs, \LcGtsrbPilotDiverged{} collapsed during training instead of failing to implant.

BPP reduces the color depth of the image rather than adding a pattern to it, and the reduced depth itself is the trigger \cite{wang2022bpp}. We quantize to \AttackBppBitDepth{} bits per channel, giving eight levels. We leave the paper's Floyd-Steinberg dithering off, since it is a sequential per-pixel operation and would slow dataset construction on Tiny ImageNet without changing which images implant. Training also uses BPP's negative samples, images quantized to a different depth drawn from 1, 2, 4, 5 or 6 bits with their label kept, at a cover rate equal to the poison rate, matching BackdoorBench's convention of one negative per poisoned image. Labeling is dirty-label all-to-one at target class 0. BPP implants at every rate on every dataset, with attack success between \AttackBppAsrMin{} and \AttackBppAsrMax{}.

Adaptive-Blend keeps Blend's whole-image pattern but poisons so that poisoned and clean images stay close in latent space, using two mechanisms together \cite{qi2023adaptive}. Cover images carry the trigger with their label kept, at a cover rate equal to the poison rate, so the latent space of a poisoned image stays close to a clean one. The pattern itself is split into a 4 by 4 grid of squares, and training reveals only half of them, chosen at random per image, while evaluation stamps the whole pattern, so the network has to generalize to a stronger trigger than it ever saw during training. Qi et al. poison at a lower opacity (\cited{0.15}) than they test at (\cited{0.2}). We use the blend ratio of Blend (\AttackAdaptiveBlendAlpha{}) for both. Labeling is dirty-label all-to-one at target class 0. This is the one attack in our set that never clears the attack success bar on any dataset or rate, reaching at most \AttackAdaptiveBlendAsrMax{}, so no Adaptive-Blend checkpoint enters the detection results at all despite being designed specifically to defeat the kind of detector PSBD is.

TaCT flips to the target only triggered images from a chosen source class and adds cover images of other classes that keep their labels, so the trigger acts only on the source class and attack success is measured on the source class alone rather than diluted by classes the attack never claims to affect \cite{tang2021demon}. The trigger is the same 3 by 3 corner checkerboard as BadNets. The source class is class 1 and the target is class 0. Training also carries cover images, the trigger stamped on non-source classes with their label kept, at a fixed cover rate of \AttackTactCoverRate{} regardless of the poison rate, since TaCT's own convention selects covers by class rather than by a multiple of the poisoning rate. TaCT implants at every rate on CIFAR-10, CIFAR-100 and Tiny ImageNet, then at 5\% and 10\% on GTSRB. Its GTSRB run at 1\% collapsed to a near-constant predictor, the same training divergence recorded for BadNets and Blend above rather than a sign that the trigger failed to implant.

Every attack is poisoned at 1\%, 5\% and 10\% of the training set, our rate ladder, except where a label mode makes a rate unreachable. The poisoned indices are drawn without replacement from a seeded generator, over the pool of images eligible under the attack's label mode. The draw is capped at the pool size when the request exceeds it. A dirty-label attack's pool is every non-target image, or every image for all-to-all, so 1\%, 5\% and 10\% are all reachable. A clean-label attack's pool is only the target class itself, so its reachable ceiling is the size of that class over the size of the training set, \CleanLabelCapCifarOneZeroClassZero{} on CIFAR-10, \CleanLabelCapCifarOneZeroZeroClassZero{} on CIFAR-100, \CleanLabelCapTinyClassZero{} on Tiny ImageNet and \CleanLabelCapGtsrbClassZero{} on GTSRB at class 0 or \CleanLabelCapGtsrbClassOne{} at class 1. A requested rate above the ceiling silently draws the same index set as the ceiling itself. A backdoored checkpoint enters the results only when two conditions hold. Its attack success rate, measured on held-out images eligible under the evaluation-time label rule rather than the training-time one, reaches \AsrBar{}. Its clean accuracy sits no further than \CleanDropBar{} below the benign reference trained with the same recipe on the same dataset. A checkpoint whose clean accuracy has collapsed below half that reference is marked diverged and excluded regardless of what its attack success reads, the rule that catches the GTSRB training collapses named above. A source-specific attack must also leave its clean source images their label. When the poison rate reaches the source class's share of the training set, every source image is poisoned and the model learns to map the whole class to the target with no trigger present. Attack success reads near 1 on such a model, but clean source images are misclassified too, so we exclude a TaCT model whose clean source-class accuracy falls below one half. That removes \PanelCellsSourceMapped{} TaCT models, every TaCT model on CIFAR-100 and Tiny ImageNet and the 10\% models on CIFAR-10 and GTSRB.

\section{From the original placement to PSBD-TM}
\label{app:placements}

\subsection{Choosing between the attention input and its twin}
\label{app:protocol}

To guard against choosing a placement on the same models it is reported on, we split the ViT models before any sweep, choosing on CIFAR-10 and GTSRB and reporting on CIFAR-100 and Tiny ImageNet. On ViT, token masking at the attention input and its twin on the attention branch output are too close for this split to separate. On the \GainsSelectionHalfN{} choosing models the twin reads \GainsSelectionHalfTwin{} against \GainsSelectionHalfOurs{}, and on the \GainsReportHalfN{} reporting models the order reverses, \GainsReportHalfOurs{} against \GainsReportHalfTwin{}. On Swin-S the same split picks the attention input by \SwinSelectionHalfTwinGap{} (\cref{sec:swin}), the attention input is also the more stable of the two across training seeds (\cref{app:seeds}) and it is the site whose mechanism \cref{sec:mechanism} explains. We therefore recommend the attention input. The numbers in the main text pool both halves.


\subsection{Dropout on the residual stream}
\label{sec:results:residual}\label{app:walk}

A ViT block has two residual additions and no activation after either, so PSBD-RD, our adaptation of the original site, applies dropout to the stream after both of them. \Cref{tab:staircase-residual} reads it beside its nearest variants, which are dropout after the attention addition only, dropout on the branch outputs just before the additions and the latter restricted to a third of the network.

% generated by scripts/paper/tab_staircase.py from results/coverage/coverage.json, results/<folder>/psbd_metrics.json (56 models) at 2026-09-30T15:01:35.277000+00:00, commit 11e403c16323e156dbac792d7a25992f49ca4dfd-dirty. Do not edit by hand.
\begin{table*}[htbp]
\centering
\small
\caption{Dropout on the residual stream of ViT-B/16, after both residual adds, after the attention add only, before both adds, and restricted to bands of blocks.}
\label{tab:staircase-residual}
\begin{tabular}{lrrrrl}
\toprule
 & & & \multicolumn{2}{c}{TPR@FPR} & \\
\cmidrule(lr){4-5}
Placement & n & AUROC & 10\% & 20\% & Gain [95\% CI] \\
\midrule
dropout, after both residual adds & 56 & 0.876 & 0.722 & 0.785 & -- \\
dropout, after the attention add only & 56 & 0.846 & 0.658 & 0.744 & $-$0.029 [$-$0.046, $-$0.015] \\
dropout, before both residual adds & 56 & 0.874 & 0.693 & 0.780 & $-$0.002 [$-$0.031, +0.027] \\
dropout, before both residual adds, blocks 1 to 4 & 56 & 0.824 & 0.579 & 0.715 & $-$0.052 [$-$0.098, $-$0.005] \\
dropout, before both residual adds, blocks 5 to 8 & 56 & 0.908 & 0.751 & 0.814 & +0.032 [+0.006, +0.062] \\
dropout, before both residual adds, blocks 9 to 12 & 41 & 0.900 & 0.776 & 0.836 & +0.006 [$-$0.042, +0.053] \\
\bottomrule
\end{tabular}
\end{table*}


PSBD-RD reaches a mean AUROC of \StaircaseResidualPostResidualAuroc{} and flags \StaircaseResidualPostResidualTprAtOneZeroPercent{} of the triggered inputs at a 10\% false-positive budget, but the mean hides a split, since it reads below chance on \StaircaseResidualPostResidualBelowChance{} of the \StaircaseResidualPostResidualN{} models. On BadNets, Blend and LF at the higher poison rates it is close to perfect, and on TaCT and on most models at \PanelRateLow{} poisoning it is close to a coin or worse. Moving the dropout to the branch outputs before the additions changes nothing on average, \StaircaseResidualPreResidualGain{} with an interval of [\StaircaseResidualPreResidualGainLow{}, \StaircaseResidualPreResidualGainHigh{}]. Restricting it to the middle or last third of the network helps a little and restricting it to the first third hurts, because the stream in early blocks carries mostly clean content.


\subsection{The same dropout at other sites}
\label{sec:results:sites}

\Cref{tab:staircase-dropout-sites} keeps the perturbation the original used and moves it to the other sites of \cref{fig:vit-block}, in every block.

% generated by scripts/paper/tab_staircase.py from results/coverage/coverage.json, results/<folder>/psbd_metrics.json (56 models) at 2026-09-30T13:16:57.522808+00:00, commit 257fa16f3d6cf6ee17ee984d58eef143db7ecc4c-dirty. Do not edit by hand.
\begin{table*}[htbp]
\centering
\small
\caption{Dropout at each site of the ViT-B/16 block, applied in all 12 blocks.}
\label{tab:staircase-dropout-sites}
\begin{tabular}{lrrrrl}
\toprule
 & & & \multicolumn{2}{c}{TPR@FPR} & \\
\cmidrule(lr){4-5}
Placement & n & AUROC & 10\% & 20\% & Gain [95\% CI] \\
\midrule
dropout, after both residual adds & 56 & 0.876 & 0.722 & 0.785 & -- \\
dropout, before both residual adds & 56 & 0.874 & 0.693 & 0.780 & $-$0.002 [$-$0.031, +0.027] \\
dropout, attention input & 56 & 0.905 & 0.708 & 0.817 & +0.029 [$-$0.021, +0.078] \\
dropout, attention input after norm & 56 & 0.846 & 0.641 & 0.749 & $-$0.030 [$-$0.081, +0.024] \\
dropout, MLP input & 56 & 0.780 & 0.543 & 0.666 & $-$0.096 [$-$0.135, $-$0.062] \\
dropout, embedding output & 56 & 0.846 & 0.579 & 0.719 & $-$0.030 [$-$0.065, +0.007] \\
\bottomrule
\end{tabular}
\end{table*}


Dropout at the attention input, before the block's LayerNorm, reaches \StaircaseDropoutSitesBeforeAttentionNormAuroc{} against \StaircaseDropoutSitesPostResidualAuroc{} on the residual stream, a paired gain of \StaircaseDropoutSitesBeforeAttentionNormGain{} [\StaircaseDropoutSitesBeforeAttentionNormGainLow{}, \StaircaseDropoutSitesBeforeAttentionNormGainHigh{}]. Dropout at the MLP input reads \StaircaseDropoutSitesBeforeMlpNormAuroc{} and dropout on the embedding output \StaircaseDropoutSitesAfterEmbeddingAuroc{}. The site matters, but the large gain comes from changing what is injected, which the main text reads in \cref{tab:staircase-operators}.


\subsection{The perturbation at matched disturbance}
\label{app:operators-matched}

\Cref{tab:staircase-operators} reads each perturbation at the rate the adaptive rule assigns it. \Cref{tab:operators-deltas} repeats the comparison with every placement read at the rate that changes the same share of clean predictions. At the attention input noise loses to token masking by \GaussianMinusTokenMaskAttentionNorm{} and to channel masking by \GaussianMinusChannelMaskAttentionNorm{}, so the ranking of the main text holds at equal disturbance. With token masking held fixed, the attention input leads the MLP input by \AttentionInputMinusMlpInputTokenMask{} [\AttentionInputMinusMlpInputTokenMaskLow{}, \AttentionInputMinusMlpInputTokenMaskHigh{}].

% generated by scripts/paper/tab_operators.py from results/coverage/coverage.json, results/<folder>/psbd_metrics.json (65 cells) at 2026-09-24T04:43:31.586866+00:00, commit 52e4ba1a7ac941365bc37d68e1e902dad4fb4f74-dirty. Do not edit by hand.
\begin{table}[htbp]
\centering
\small
\caption{Paired AUROC differences between placements on the same ViT-B/16 models at matched disturbance, a clean shift ratio of 0.6, with 5000-resample bootstrap 95\% intervals. The first three rows change the perturbation at a fixed site and the last changes the site at a fixed perturbation.}
\label{tab:operators-deltas}
\begin{adjustbox}{max width=\linewidth}
\begin{tabular}{lrrl}
\toprule
comparison & models & mean difference & 95\% interval \\
\midrule
noise minus token mask, attention input & 65 & $-$0.206 & [$-$0.259, $-$0.156] \\
noise minus channel mask, attention input & 65 & $-$0.113 & [$-$0.151, $-$0.078] \\
noise after the MLP norm minus token mask before it & 65 & +0.056 & [+0.026, +0.087] \\
attention input minus MLP input, both token mask & 65 & +0.109 & [+0.079, +0.138] \\
\bottomrule
\end{tabular}
\end{adjustbox}
\end{table}

\FloatBarrier


\subsection{Masking in every block}
\label{app:bands}

\Cref{tab:depth-bands} restricts token masking at the attention input to one third of the network (blocks 1 to 4, 5 to 8 or 9 to 12) and reads each band at matched disturbance. Every band loses to masking in every block, and the first third loses most, \BandOneFourMinusAllInputSide{} [\BandOneFourMinusAllInputSideLow{}, \BandOneFourMinusAllInputSideHigh{}], as expected for a backdoor written late in the network.

% generated by scripts/paper/tab_depth_bands.py from results/coverage/coverage.json, results/<folder>/psbd_metrics.json (56 cells) at 2026-09-29T18:04:10.632150+00:00, commit b2d32cf11708d5de965d3e13604c863a3ad9b493-dirty. Do not edit by hand.
\begin{table}[htbp]
\centering
\small
\caption{Depth-band placements at the matched 0.6 rate, mean AUROC at the headline quantile, the paired delta against the family's all-blocks placement with a 5000-resample bootstrap 95\% interval, and the mean achieved clean-validation shift ratio at the chosen rate. The n column is the models behind each row, which differs by band because the bands were swept at different times.}
\label{tab:depth-bands}
\begin{adjustbox}{max width=\linewidth}
\begin{tabular}{lrrrlr}
\toprule
placement & n & mean AUROC matched06 & delta vs all blocks & 95\% CI & mean achieved shift \\
\midrule
input-side, all blocks & 54 & 0.938 & -- & -- & 0.594 \\
input-side, blocks 5-8 & 54 & 0.867 & $-$0.071 & [$-$0.112, $-$0.034] & 0.603 \\
input-side, blocks 9-12 & 54 & 0.886 & $-$0.052 & [$-$0.106, $-$0.006] & 0.482 \\
input-side, blocks 1-4 & 54 & 0.833 & $-$0.104 & [$-$0.138, $-$0.073] & 0.534 \\
residual-adjacent, all blocks & 54 & 0.863 & -- & -- & 0.667 \\
residual-adjacent, blocks 1-4 & 54 & 0.786 & $-$0.077 & [$-$0.097, $-$0.057] & 0.573 \\
residual-adjacent, blocks 5-8 & 54 & 0.911 & +0.048 & [+0.021, +0.076] & 0.562 \\
residual-adjacent, blocks 9-12 & 54 & 0.899 & +0.036 & [$-$0.009, +0.080] & 0.596 \\
\bottomrule
\end{tabular}
\end{adjustbox}
\end{table}

\FloatBarrier


\subsection{Other perturbations of attention}
\label{app:attention}

Attention is the only operation in a ViT block that moves information between tokens, so we measured \AttentionProbesTotal{} perturbations of it in all. \Cref{tab:attention-probes} gives them. Only the twin on the attention branch output ties with token masking at the attention input, and none beats it. Masking whole attention heads, the most attention-specific of them, reads \AttentionHeadMaskAuroc{} on \AttentionHeadMaskN{} models.

% generated by scripts/paper/tab_attention.py from results/coverage/coverage.json (56 successful cells), results/*/psbd_metrics.json at 2026-09-29T18:02:53.295782+00:00, commit b2d32cf11708d5de965d3e13604c863a3ad9b493-dirty. Do not edit by hand.
\begin{table}[htbp]
\centering
\small
\caption{Every attention-side probe with at least 5 backdoored ViT-B/16 models, at the adaptive rule and the headline quantile, ordered by mean AUROC. The number of models differs between rows, so the rows are not paired.}
\label{tab:attention-probes}
\begin{adjustbox}{max width=\linewidth}
\begin{tabular}{llrrrr}
\toprule
probe & n & AUROC & TPR@10 & TPR@20 \\
\midrule
token mask, attention input & 54 & 0.963 & 0.89 & 0.92 \\
token mask, attention output before the add & 54 & 0.956 & 0.86 & 0.92 \\
channel mask, attention input & 54 & 0.921 & 0.77 & 0.84 \\
noise, attention input after norm & 19 & 0.913 & 0.79 & 0.87 \\
dropout, attention input & 54 & 0.903 & 0.71 & 0.82 \\
head mask, attention heads & 36 & 0.899 & 0.75 & 0.85 \\
token mask, attention input after norm & 36 & 0.877 & 0.75 & 0.81 \\
drop path, attention output before the add & 36 & 0.875 & 0.67 & 0.79 \\
dropout, stream after the attention add & 54 & 0.856 & 0.67 & 0.76 \\
channel mask, attention input after norm & 36 & 0.849 & 0.67 & 0.77 \\
dropout, attention input after norm & 54 & 0.847 & 0.64 & 0.75 \\
dropout, attention output before the add & 54 & 0.843 & 0.64 & 0.74 \\
noise, attention output before the add & 54 & 0.840 & 0.65 & 0.73 \\
noise, attention input & 54 & 0.836 & 0.60 & 0.71 \\
channel mask, stream after the attention add & 36 & 0.833 & 0.65 & 0.73 \\
noise, stream after the attention add & 36 & 0.820 & 0.59 & 0.70 \\
channel mask, attention output before the add & 54 & 0.817 & 0.61 & 0.71 \\
token mask, stream after the attention add & 54 & 0.805 & 0.53 & 0.63 \\
gain scale, attention norm output & 36 & 0.638 & 0.39 & 0.46 \\
\bottomrule
\end{tabular}
\end{adjustbox}
\end{table}

\FloatBarrier


\subsection{The shift ratio target}
\label{app:rate}

The original rule reads each placement at the rate where a share \AdaptiveShiftTarget{} of clean validation predictions change. We checked that target by interpolating every model's rate ladder onto a shared grid of achieved shift ratios and reading AUROC and TPR along it. AUROC barely moves with the target, from \LadderAurocAtZeroSix{} at a shift ratio of \PlacementMatchTarget{} to \LadderAurocAtZeroEight{} at \AdaptiveShiftTarget{}, while the true-positive rate at a 10\% budget rises from \LadderTprOneZeroAtZeroSix{} to \LadderTprOneZeroAtZeroEight{}. Only \LadderAurocAtZeroNineN{} of the \LadderAurocAtZeroEightN{} models reach a shift ratio of \LadderStretchTarget{} within the swept rates, so we keep the original target.


\FloatBarrier
\clearpage

\section{The mechanism in detail}
\label{app:mechanism-group}

\subsection{The depth profile and the routing, per block}
\label{app:mechanism}

% generated by scripts/paper/mech_direction.py from results/_experiments/backdoor_neurons/backdoor_neuron_ablation.json at 2026-09-30T13:02:27.253380+00:00, commit 257fa16f3d6cf6ee17ee984d58eef143db7ecc4c-dirty. Do not edit by hand.
\begin{table}[htbp]
\centering
\small
\caption{Attack success on CIFAR-10 at the highest panel rate before any ablation and after projecting out, at the peak layer (block 12), the rank-1 backdoor direction, a random direction of the same norm, the 20 coordinates the trigger moves most and 20 random coordinates. Clean accuracy falls by at most 0.038 under any of them.}
\label{tab:direction-ablation}
\begin{adjustbox}{max width=\linewidth}
\begin{tabular}{lrrrrr}
\toprule
attack & none & direction & random dir. & top 20 & random 20 \\
\midrule
BadNets & 1.000 & 0.000 & 1.000 & 1.000 & 1.000 \\
Blend & 1.000 & 0.026 & 1.000 & 1.000 & 1.000 \\
BPP & 0.998 & 0.085 & 0.998 & 0.998 & 0.998 \\
LF & 1.000 & 0.572 & 1.000 & 1.000 & 1.000 \\
\bottomrule
\end{tabular}
\end{adjustbox}
\end{table}


\begin{figure*}[t]
\centering
\includegraphics[width=\textwidth]{mech_activation_patching.pdf}
\caption{Share of the clean prediction recovered when a token group is replaced by its clean value at a given block, averaged over the backdoored models of each attack. For the local triggers (BadNets, TaCT) the trigger's tokens hold the backdoor through block 10 and the class token takes over only at the end.}
\label{fig:patching}
\end{figure*}

The direction is not present from the start. \Cref{fig:tac-layers} follows its norm, relative to the clean activation's norm, through the \VitBlocks{} blocks on GTSRB, CIFAR-100 and Tiny ImageNet. On the benign model it never exceeds \DirectionBenignPeakRelNorm{}. On a backdoored model it stays near the benign level through the first half of the network and reaches half of its final value between blocks 7 and 9 depending on the attack. It peaks at blocks 9 to 12. The onset comes earlier for the global triggers than for the patch trigger on CIFAR-100 and Tiny ImageNet, and the order reverses on GTSRB, so the depth at which a trigger is written depends on the dataset as well as the trigger.

\begin{figure*}[t]
\centering
\includegraphics[width=\textwidth]{mech_tac_layers.pdf}
\caption{Relative norm of the backdoor direction (top) and mean trigger-activated change (TAC) \cite{zheng2022clp} per dimension, the mean absolute activation difference between a triggered image and its clean copy (bottom), by block, on GTSRB, CIFAR-100 and Tiny ImageNet at 10\% poisoning, \TacLayersSamples{} paired images per model. The benign model probed with the same trigger is dashed. The backdoor is written in the second half of the network.}
\label{fig:tac-layers}
\end{figure*}

The class token's attention confirms the route. On the BadNets models the share of its attention that falls on the trigger's four tokens is \RoutingBadnetClsEarly{} in blocks 1 to 4, \RoutingBadnetClsRouting{} in blocks 5 to 8 and \RoutingBadnetClsLate{} in blocks 9 to 12, against \RoutingBenignClsLate{} on the benign model probed with the same trigger, over \RoutingCells{} checkpoints. Attention, the only operation that moves information between tokens, carries the trigger to the class token in the second half of the network. \Cref{fig:routing} shows this per block.

\begin{figure*}[t]
\centering
\includegraphics[width=\textwidth]{mech_routing.pdf}
\caption{Left, share of the class token's attention that falls on the trigger's tokens, by block and attack, against the benign model probed with the same trigger (dashed). Right, the norm rank of the trigger's tokens among all patch tokens. The trigger is routed to the class token in blocks 5 to 12.}
\label{fig:routing}
\end{figure*}


\subsection{Which predictions survive, per attack}
\label{app:survival}

\Cref{tab:survival} gives, per attack, the share of triggered predictions that each placement changes at its adaptive rate and the resulting AUROC, the numbers behind \cref{fig:survival}. \Cref{tab:survival-sites} reads BadNets under each perturbation of \cref{sec:mechanism:placements}, and \cref{tab:causal-patch} gives the causal test, perturbing only the trigger's tokens or as many random ones.

% generated by scripts/paper/mech_survival.py from results/coverage/coverage.json, results/<folder>/psbd_metrics.json (57 models) at 2026-09-24T06:09:59.984864+00:00, commit a6bece9dac4f610c3f14c6d149a3dd576da74d71-dirty. Do not edit by hand.
\begin{table}[htbp]
\centering
\small
\caption{Share of triggered predictions that change under each placement at its adaptive rate, where about four in five clean predictions change, and the resulting mean AUROC. Patch triggers first, then global triggers. A detector needs the triggered share far below the clean one.}
\label{tab:survival}
\begin{adjustbox}{max width=\linewidth}
\begin{tabular}{lrrrrr}
\toprule
attack & n & TM shift & TM AUROC & RD shift & RD AUROC \\
\midrule
BadNets & 12 & 0.04 & 0.992 & 0.61 & 0.712 \\
TaCT & 3 & 0.80 & 0.962 & 0.94 & 0.563 \\
Blend & 12 & 0.04 & 0.978 & 0.14 & 0.971 \\
SIG & 1 & 0.90 & 0.418 & 0.25 & 0.919 \\
WaNet & 5 & 0.23 & 0.845 & 0.11 & 0.956 \\
LF & 12 & 0.03 & 0.980 & 0.14 & 0.974 \\
BPP & 12 & 0.10 & 0.948 & 0.20 & 0.945 \\
\bottomrule
\end{tabular}
\end{adjustbox}
\end{table}

% generated by scripts/paper/mech_survival.py from results/coverage/coverage.json, results/<folder>/psbd_metrics.json (56 models) at 2026-09-30T14:51:27.149631+00:00, commit 1fd38381a58031e98490ed24679d373b446725fb-dirty. Do not edit by hand.
\begin{table}[htbp]
\centering
\small
\caption{BadNets under each perturbation at its adaptive rate. The clean column is the share of clean test predictions that change, the triggered column the share of triggered predictions that change, and AUROC the resulting detection.}
\label{tab:survival-sites}
\begin{adjustbox}{max width=\linewidth}
\begin{tabular}{lrrrr}
\toprule
perturbation & n & clean & triggered & AUROC \\
\midrule
token mask, attention input & 12 & 0.87 & 0.04 & 0.992 \\
token mask, attention branch output & 12 & 0.90 & 0.10 & 0.973 \\
token mask, stream after the attention add & 12 & 0.87 & 0.58 & 0.812 \\
dropout, attention input & 12 & 0.83 & 0.33 & 0.865 \\
channel mask, attention input & 12 & 0.88 & 0.37 & 0.881 \\
Gaussian noise, attention input & 12 & 0.87 & 0.64 & 0.626 \\
dropout, stream after both adds (PSBD-RD) & 12 & 0.88 & 0.61 & 0.712 \\
\bottomrule
\end{tabular}
\end{adjustbox}
\end{table}

% generated by scripts/paper/mech_causal.py from results/_experiments/why_token_masking_works/summary.json (39 models) at 2026-09-29T16:15:58.883088+00:00, commit 2e59453ac935840c88040509aab36fb7e2c4ac3d-dirty. Do not edit by hand.
\begin{table}[htbp]
\centering
\small
\caption{Share of predictions kept on the 12 BadNets models, triggered and clean, when only chosen tokens are perturbed. The masks are deterministic. The dropout rows use PSBD-RD's adaptive rate with ten passes.}
\label{tab:causal-patch}
\begin{adjustbox}{max width=\linewidth}
\begin{tabular}{lrr}
\toprule
perturbed & triggered kept & clean kept \\
\midrule
\multicolumn{3}{l}{\emph{masked at the attention input}} \\
trigger tokens, all 12 blocks & 0.002 & 0.986 \\
trigger tokens, blocks 1 to 4 & 1.000 & 0.988 \\
trigger tokens, blocks 5 to 8 & 0.788 & 0.995 \\
trigger tokens, blocks 9 to 12 & 0.200 & 0.991 \\
random tokens, all 12 blocks & 1.000 & 0.983 \\
\multicolumn{3}{l}{\emph{PSBD-RD dropout applied to}} \\
every position (PSBD-RD) & 0.398 & 0.113 \\
trigger positions only & 0.387 & 0.995 \\
every position but the trigger & 0.664 & 0.119 \\
as many random positions & 1.000 & 0.992 \\
\bottomrule
\end{tabular}
\end{adjustbox}
\end{table}

\FloatBarrier


\subsection{LayerNorm absorbs noise and not masking}
\label{app:layernorm}\label{sec:mechanism:layernorm}

A LayerNorm rescales each token to unit variance, so additive noise that raises the variance is partly divided away and a mask is not. The attention input's LayerNorm passes \AbsorptionAttentionInputGaussianSurvival{} of a Gaussian disturbance and \AbsorptionAttentionInputTokenMaskSurvival{} of a token mask on \AbsorptionModels{} models, against \AbsorptionControlSurvival{} where no LayerNorm follows. The adaptive rule therefore picks a larger rate for noise, and noise injected after the MLP's LayerNorm escapes the absorption and reads \StaircaseOperatorsBeforeMlpGaussianAuroc{}, the apparent reversal of \cref{sec:results:operator}.

\Cref{tab:layernorm-absorption} gives the share of each perturbation that survives the LayerNorm behind its site.

% generated by scripts/paper/mech_layernorm.py from results/*/layernorm_absorption.json (2 models) at 2026-09-24T04:31:01.535219+00:00, commit 52e4ba1a7ac941365bc37d68e1e902dad4fb4f74-dirty. Do not edit by hand.
\begin{table}[htbp]
\centering
\small
\caption{What a LayerNorm passes through. Survival is the perturbation's relative norm after the LayerNorm over its relative norm before it, averaged over blocks and over two models. The last two rows sit after the LayerNorm, where nothing normalizes them, and are the control.}
\label{tab:layernorm-absorption}
\begin{adjustbox}{max width=\linewidth}
\begin{tabular}{llrrr}
\toprule
position & a norm follows & token mask & channel mask & gaussian \\
\midrule
attention input & yes & 1.002 & 0.968 & 0.821 \\
MLP input & yes & 0.971 & 0.981 & 0.872 \\
attention input after norm & no & 1.000 & 1.000 & 1.000 \\
MLP input after norm & no & 1.000 & 1.000 & 1.000 \\
\bottomrule
\end{tabular}
\end{adjustbox}
\end{table}

\FloatBarrier


\subsection{The original account holds for some attacks}
\label{app:shift-target}\label{sec:mechanism:phenomenon}

The PSBD paper explains its statistic by neuron bias, under which clean inputs whose prediction shifts should land on the backdoor's target. On ViT \ShiftToTargetDriftAttacks{} drift toward the target on some dataset, while \ShiftToTargetNoDriftAttacks{} stay below \ShiftToTargetNoDriftFloor{} on every dataset. PSBD-TM detects Blend and TaCT, which show no drift, so neuron bias describes some attacks rather than the statistic's mechanism.

Among clean validation images whose prediction changes under PSBD-TM, the mean excess of the target share over the benign reference is \ShiftToTargetExcessAll{}, and it is concentrated in a few attacks, the largest excess \ShiftToTargetLargestExcess{} for \ShiftToTargetLargestCell{}. The split does not follow the kind of trigger, since TaCT is a local patch and does not drift while LF and BPP are global and do. Nor is the statistic the unperturbed confidence. The maximum softmax probability reads a mean AUROC of \ConfidenceNullAuroc{} over \ConfidenceNullCheckpoints{} models against \ConfidenceNullPsuAuroc{} for PSU, and PSU is ahead on \ConfidenceNullPsuWins{} of them.

\Cref{tab:mech-shift-target-by-dataset} gives the share of shifted clean predictions that land on the target class, per dataset and attack, against the benign model of the same dataset.


% generated by scripts/paper/mech_shift_target.py from results/<folder>/psbd_metrics.json (the successful cells and the benign references) at 2026-09-30T14:50:35.252650+00:00, commit 1fd38381a58031e98490ed24679d373b446725fb-dirty. Do not edit by hand.
\begin{table}[htbp]
\centering
\small
\caption{Share of shifted clean predictions that land on the target class under PSBD-TM, per dataset and attack, against the uniform expectation and the benign reference of the dataset.}
\label{tab:mech-shift-target-by-dataset}
\begin{adjustbox}{max width=\linewidth}
\begin{tabular}{lrrrrrr}
\toprule
Dataset & Attack & n & Share & Uniform & Benign & Excess \\
\midrule
CIFAR-10 & BPP & 3 & 0.166 & 0.100 & 0.004 & +0.162 \\
CIFAR-10 & WaNet & 1 & 0.002 & 0.100 & 0.004 & $-$0.002 \\
CIFAR-10 & TaCT & 2 & 0.000 & 0.100 & 0.004 & $-$0.004 \\
CIFAR-10 & BadNets & 3 & 0.518 & 0.100 & 0.004 & +0.514 \\
CIFAR-10 & Blend & 3 & 0.005 & 0.100 & 0.004 & +0.001 \\
CIFAR-10 & LF & 3 & 0.172 & 0.100 & 0.004 & +0.168 \\
CIFAR-100 & BPP & 3 & 0.157 & 0.010 & 0.000 & +0.157 \\
CIFAR-100 & BadNets & 3 & 0.002 & 0.010 & 0.000 & +0.002 \\
CIFAR-100 & Blend & 3 & 0.009 & 0.010 & 0.000 & +0.009 \\
CIFAR-100 & LF & 3 & 0.256 & 0.010 & 0.000 & +0.256 \\
GTSRB & BPP & 3 & 0.001 & 0.023 & 0.000 & +0.001 \\
GTSRB & TaCT & 2 & 0.000 & 0.023 & 0.000 & +0.000 \\
GTSRB & Label-Consistent & 1 & 0.000 & 0.023 & 0.000 & +0.000 \\
GTSRB & BadNets & 3 & 0.793 & 0.023 & 0.000 & +0.793 \\
GTSRB & Blend & 3 & 0.015 & 0.023 & 0.000 & +0.015 \\
GTSRB & LF & 3 & 0.054 & 0.023 & 0.000 & +0.054 \\
Tiny ImageNet & BPP & 3 & 0.347 & 0.005 & 0.002 & +0.346 \\
Tiny ImageNet & WaNet & 2 & 0.400 & 0.005 & 0.002 & +0.398 \\
Tiny ImageNet & BadNets & 3 & 0.150 & 0.005 & 0.002 & +0.149 \\
Tiny ImageNet & Blend & 3 & 0.005 & 0.005 & 0.002 & +0.003 \\
Tiny ImageNet & LF & 3 & 0.385 & 0.005 & 0.002 & +0.383 \\
\bottomrule
\end{tabular}
\end{adjustbox}
\end{table}

\FloatBarrier


\section{The adaptive attacker, the probe union and forward passes}
\label{app:robustness-group}

\subsection{The adaptive attacker}
\label{app:adaptive}

\Cref{tab:adaptive} summarizes the evasive models behind \cref{sec:robustness:adaptive}, and \cref{tab:adaptive-attacker} in the supplement gives each of them.

% generated by scripts/paper/tab_adaptive.py from results/adaptive_attacker_analysis.json, results/multi_probe_analysis.json at 2026-09-24T05:46:23.024430+00:00, commit 872fbbbd1f025c8c9924211876e2cda72fc45aab-dirty. Do not edit by hand.
\begin{table}[htbp]
\centering
\small
\caption{The adaptive attacker on the CIFAR-100 ViT-B/16 models whose evasive version keeps attack success at or above 0.9. ASR is the attack success rate of the evasive model and cost its clean-accuracy loss. Probed is the AUROC of the placement the attacker trained against, transfer the mean AUROC of the two probes the attacker never saw and union the AUROC of the min-rank union of all three. \Cref{tab:adaptive-attacker} gives every model separately.}
\label{tab:adaptive}
\begin{adjustbox}{max width=\linewidth}
\begin{tabular}{lrrrrr}
\toprule
 & ASR & cost & probed & transfer & union \\
\midrule
before the attack & 0.966 & -- & 0.966 & -- & -- \\
after the attack, mean of 14 & 0.981 & $-$0.030 & \textbf{0.233} & 0.893 & \textbf{0.967} \\
worst model for the probe (BadNets 5\%) & 1.000 & $-$0.026 & 0.004 & 0.785 & 0.994 \\
worst model for the union (LF 1\%) & 0.938 & $-$0.042 & 0.410 & 0.838 & 0.852 \\
\bottomrule
\end{tabular}
\end{adjustbox}
\end{table}

\FloatBarrier


\subsection{The probe union on the ordinary models}
\label{app:union}

The union is not only for the adaptive case. \Cref{tab:probe-union} reads the same rule on the \ProbeUnionNModels{} ordinary models, where nobody trained against a probe. Adding PSBD-RD to PSBD-TM rescues the two models PSBD-TM misses, WaNet on CIFAR-10 from \ProbeUnionWanetCifarOneZeroTm{} to \ProbeUnionWanetCifarOneZeroTmRd{} and SIG on CIFAR-10 from \HeadlineFloorAuroc{} to a comparable level. It costs on TaCT, where PSBD-RD reads a coin and its rank drags the minimum down, so the paired gain over all models is \ProbeUnionTmRdGain{} with an interval \ProbeUnionTmRdGainCi{} that crosses zero. Pairing PSBD-TM with token masking on the attention branch output instead gains \ProbeUnionTmBranchGain{} with an interval that excludes zero, because the two probes fail on different models. That is the general case rather than luck, since placements disagree on which inputs are fragile. Per-sample PSU under token masking at the attention input correlates at \AgreementGaussianVsTokenMaskClean{} with noise at the same site and at \AgreementGainScaleVsRecommendedClean{} with gain scaling of the MLP LayerNorm output (\cref{tab:mech-operator-agreement}), so a second probe carries information the first does not. The three-probe union of the adaptive section reads \ProbeUnionThreeProbeAuroc{} on the \ProbeUnionThreeProbeN{} models holding all three probes, and a union of every fully swept placement gains \ProbeUnionAllBasisGain{}, about what the right pair gains. More probes help when they disagree in the right way, not by number.

\IfFileExists{tables/probe_union.tex}{% generated by scripts/paper/tab_probe_union.py from results/_experiments/probe_union/probe_union.json at 2026-09-24T06:07:06.287836+00:00, commit a6bece9dac4f610c3f14c6d149a3dd576da74d71-dirty. Do not edit by hand.
\begin{table*}[htbp]
\centering
\small
\caption{The min-rank union of probes on the 57 backdoored ViT-B/16 models, with the paired AUROC gain over token masking at the attention input alone. TM is token masking at the attention input, RD dropout after both residual adds, TM-out token masking on the attention branch output, DI dropout at the attention input and GS gain scaling of the MLP LayerNorm output.}
\label{tab:probe-union}
\begin{adjustbox}{max width=\textwidth}
\begin{tabular}{lrrrrl}
\toprule
Probes & n & AUROC & TPR@FPR 10\% & TPR@FPR 20\% & Gain over TM [95\% CI] \\
\midrule
TM alone & 57 & 0.953 & 0.873 & 0.902 & -- \\
TM + RD & 57 & 0.969 & 0.900 & 0.928 & +0.017 [$-$0.005, +0.044] \\
TM + TM-out & 57 & 0.973 & 0.908 & 0.941 & +0.021 [+0.003, +0.045] \\
TM + DI + GS & 57 & 0.967 & 0.907 & 0.934 & +0.015 [$-$0.002, +0.035] \\
TM + DI + GS + RD & 57 & 0.975 & 0.919 & 0.948 & +0.022 [+0.000, +0.049] \\
every basis placement on every model & 57 & 0.976 & 0.913 & 0.943 & +0.024 [+0.002, +0.052] \\
\bottomrule
\end{tabular}
\end{adjustbox}
\end{table*}
}{}


\subsection{Forward passes and training seeds}
\label{app:cost}

\paragraph{Forward passes.}
\label{sec:robustness:passes}

PSBD uses $k = 3$ passes. Over all \PassAllCells{} models the mean AUROC of PSBD-TM rises from \PassAllAurocKOne{} at $k = 1$ to \PassAllAurocKTwo{} at $k = 2$ and \PassAllAurocKThree{} at $k = 3$. On the \PassPilotCells{} models whose sweep reached 20 passes it rises from \PassPilotAurocKOne{} to \PassPilotAurocKThree{} at three passes and \PassPilotAurocKTwoZero{} at 20, with TPR at a 10\% false-positive budget going from \PassPilotTprOneZeroKOne{} to \PassPilotTprOneZeroKThree{} to \PassPilotTprOneZeroKTwoZero{}. Beyond the third pass the gain is \PassPilotAurocGainBeyondThree{} AUROC, so we keep the original $k = 3$. \Cref{app:passes} plots both curves.

\paragraph{Seed variance.}
\label{sec:robustness:seeds}

We trained \SeedCells{} backdoored models \SeedsPerCell{} times with different seeds (BadNets, Blend and WaNet on CIFAR-100 and Tiny ImageNet). PSBD-TM has a seed standard deviation of \SeedSdRecommended{} AUROC on average and PSBD-RD one of \SeedSdPublished{}. The gain of PSBD-TM over PSBD-RD is \SeedGainMean{} on these models and changes sign on \SeedGainSignFlips{} of \SeedCells{}, all on Blend and WaNet where both placements are near the ceiling. A single seed is therefore a fair reading of PSBD-TM and an unreliable reading of PSBD-RD. \Cref{app:seeds} lists the per-seed readings.

\subsection{The forward-pass curve}
\label{app:passes}

\Cref{fig:forward-passes} plots the effect of the number of passes $k$ behind \cref{sec:robustness:passes}. Over all models the first three passes are available from the sweeps, and on \PassPilotCells{} models at 1\% poisoning we swept $k$ up to 20.

\begin{figure}[htbp]
\centering
\includegraphics[width=\textwidth]{fig_forward_passes.pdf}
\caption{Mean AUROC, TPR at 10\% FPR and TPR at 20\% FPR of PSBD-TM against the number of forward passes $k$. The full curve is on the \PassPilotCells{} models at 1\% poisoning with a $k = 20$ sweep. The short curve is on all models, whose sweeps hold $k \le 3$.}
\label{fig:forward-passes}
\end{figure}


\subsection{Seed replicates}
\label{app:seeds}

\Cref{tab:seeds} gives the readings behind \cref{sec:robustness:seeds}. Each row is a backdoored model trained \SeedsPerCell{} times with different seeds. The columns give the smallest attack success rate among the runs, the AUROC of PSBD-TM on each seed with its standard deviation, the same for PSBD-RD and the mean and standard deviation of their paired gain. PSBD-TM's seed standard deviation averages \SeedSdRecommended{} against \SeedSdPublished{} for PSBD-RD, and the paired gain's own standard deviation reaches \SeedSdGainMax{} on its worst row.

% generated by scripts/paper/tab_seeds.py from results/coverage/coverage.json, results/<folder>[_seed_N]/psbd_metrics.json (14 cells), checkpoints/<folder>_seed_N/args.json at 2026-09-24T05:55:23.343403+00:00, commit 872fbbbd1f025c8c9924211876e2cda72fc45aab-dirty. Do not edit by hand.
\begin{table}[htbp]
\centering
\small
\caption{PSBD-TM and PSBD-RD on the models trained with three seeds, as the mean AUROC over seeds and its standard deviation. Min ASR is the lowest attack success among the runs of a model.}
\label{tab:seeds}
\begin{adjustbox}{max width=\linewidth}
\begin{tabular}{lllrrrrrrr}
\toprule
Dataset & Attack & Rate & Min ASR & PSBD-TM & sd & PSBD-RD & sd & gain & sd \\
\midrule
CIFAR-100 & BadNets & 0.01 & 0.999 & 0.989 & 0.006 & 0.572 & 0.151 & +0.417 & 0.151 \\
CIFAR-100 & BadNets & 0.05 & 1.000 & 0.980 & 0.026 & 0.531 & 0.278 & +0.449 & 0.259 \\
CIFAR-100 & BadNets & 0.1 & 1.000 & 0.997 & 0.002 & 0.574 & 0.081 & +0.423 & 0.081 \\
CIFAR-100 & Blend & 0.01 & 0.988 & 0.963 & 0.036 & 0.885 & 0.100 & +0.079 & 0.095 \\
CIFAR-100 & Blend & 0.05 & 0.999 & 0.986 & 0.009 & 0.966 & 0.055 & +0.020 & 0.053 \\
CIFAR-100 & Blend & 0.1 & 0.999 & 0.994 & 0.005 & 0.996 & 0.004 & $-$0.001 & 0.003 \\
Tiny ImageNet & BadNets & 0.01 & 1.000 & 0.980 & 0.004 & 0.876 & 0.042 & +0.104 & 0.038 \\
Tiny ImageNet & BadNets & 0.05 & 1.000 & 0.984 & 0.008 & 0.813 & 0.107 & +0.171 & 0.099 \\
Tiny ImageNet & BadNets & 0.1 & 1.000 & 0.983 & 0.012 & 0.827 & 0.164 & +0.156 & 0.151 \\
Tiny ImageNet & Blend & 0.01 & 1.000 & 0.996 & 0.003 & 0.997 & 0.001 & $-$0.001 & 0.002 \\
Tiny ImageNet & Blend & 0.05 & 0.997 & 0.991 & 0.011 & 0.998 & 0.003 & $-$0.006 & 0.007 \\
Tiny ImageNet & Blend & 0.1 & 1.000 & 0.997 & 0.001 & 1.000 & 0.000 & $-$0.002 & 0.001 \\
Tiny ImageNet & WaNet & 0.05 & 0.828 & 0.918 & 0.039 & 0.931 & 0.047 & $-$0.013 & 0.011 \\
Tiny ImageNet & WaNet & 0.1 & 0.931 & 0.955 & 0.010 & 0.956 & 0.024 & $-$0.001 & 0.018 \\
\bottomrule
\end{tabular}
\end{adjustbox}
\end{table}

\FloatBarrier


\FloatBarrier
\clearpage

\section{Swin-S in detail}
\label{app:swin-group}

\subsection{Swin-S paired comparisons}
\label{app:swin}

\Cref{tab:swin-gains} gives the paired comparisons behind \cref{sec:swin} with their bootstrap intervals. \Cref{tab:swin-attacks} in the supplement gives every Swin-S model separately and \cref{tab:swin-placements} every placement cached on Swin with the number of models behind each.

% generated by scripts/paper/tab_swin.py from configs/psbd_basis.json (panel rule retargeted to swin, 93 panel cells), results/swin_*/psbd_metrics.json (65 successful cells), checkpoints/swin_*/args.json at 2026-09-29T18:13:30.484532+00:00, commit d73d3277d007fd75b3a4166a7911b1e07a2761b7-dirty. Do not edit by hand.
\begin{table}[htbp]
\centering
\small
\caption{Swin-S paired placement gains at the headline quantile over the models carrying both placements, 5000-resample bootstrap intervals.}
\label{tab:swin-gains}
\begin{adjustbox}{max width=\linewidth}
\begin{tabular}{llrrl}
\toprule
comparison & rule & n & mean delta AUROC & 95\% CI \\
\midrule
token mask, attention input minus dropout, after both residual adds & matched & 65 & +0.135 & [+0.089, +0.187] \\
token mask, attention input minus dropout, after both residual adds & adaptive & 63 & +0.096 & [+0.051, +0.144] \\
dropout, attention input minus dropout, after both residual adds & matched & 63 & +0.009 & [$-$0.039, +0.056] \\
dropout, attention input minus dropout, after both residual adds & adaptive & 63 & +0.015 & [$-$0.046, +0.074] \\
token mask, attention input minus dropout, before both residual adds & matched & 63 & +0.140 & [+0.092, +0.191] \\
token mask, attention input minus dropout, before both residual adds & adaptive & 63 & +0.112 & [+0.059, +0.169] \\
token mask, attention input minus token mask, attention output before the add & matched & 65 & +0.151 & [+0.105, +0.202] \\
token mask, attention input minus token mask, attention output before the add & adaptive & 63 & +0.128 & [+0.070, +0.187] \\
token mask, attention input minus dropout, attention input & matched & 63 & +0.128 & [+0.084, +0.176] \\
token mask, attention input minus dropout, attention input & adaptive & 63 & +0.081 & [+0.042, +0.127] \\
\bottomrule
\end{tabular}
\end{adjustbox}
\end{table}


\FloatBarrier
\clearpage
